
\chapter{Rappels d'analyse vectorielle}
% \section*{Rappels d'analyse vectorielle}
Soient un vecteur $P$ et un réel $\lambda$. Rappelons que 
\begin{equation}
\label{eq:div_scal}
\nabla\cdot(\lambda P) = P\cdot\nabla\lambda\ + \lambda\nabla\cdot P.
\end{equation}


Rappelons aussi la relation bien connue :
\begin{equation}
\label{eq:vect_id}
\helm P = \nabla\left(\nabla\cdot P\right)-\nabla^2P.
\end{equation}

\'Etant donné un vecteur unitaire $n$ normal au vecteur $P$, on remarque que :
\begin{equation}
	\label{eq:vect_normale}
	n\times P\times n = P.
\end{equation}

Soit $ \Gamma$ une surface gauche dans $\mathbb R^3$. Soit une fonction $v$ régulière définie au voisinage de $\Gamma$. On peut définir un gradient surfacique $\nabla_\Gamma$ de $v|_\Gamma$ comme la composante tangentielle à $\Gamma$ de $\nabla v$ :
\begin{equation}
    \left.\nabla_\Gamma v\right|_\Gamma = \n\times\left(\nabla v\times\n\right).
    \label{eq:gradsurf}
\end{equation}

\'Etant donné un champ de vecteur $\J$ régulier et tangent à $\Gamma$, on peut alors définir sa divergence surfacique $\nabla_\Gamma\cdot$ par dualité :
\begin{equation}
    \int_\Gamma\phi\nabla_\Gamma\cdot\J\dd\Gamma = - \int_\Gamma\J\cdot\nabla_\Gamma\phi\dd\Gamma
    \label{eq:divsurf}
\end{equation}

Soient $u$, $v$ et $w$ trois vecteurs de $\mathbb R^3$. Le produit vectoriel n'est pas associatif, et :
\begin{align}
    u\times(v\times w) &= (u\cdot w)v - (u\cdot v)w\\
    (u\times v)\times w &= (u\cdot w)v - (v\cdot w)u.
\end{align}

Cependant, si $w = u$, on a :
\begin{equation}
    u\times(v\times u) = (u\cdot u)v - (u\cdot v)u = (u\times v)\times u.
\end{equation}

En particulier, si $u\perp v$, alors :
\begin{equation}
    u\times v\times u = (u\cdot u)v - (v\cdot u)u = |u|^2 v
    \label{eq:prodvec}
\end{equation}